% GENERATED FILE. Edit canonical lessons, then run npm run generate:books.
% canonical-source-hash: fnv1a64:443bd7f8dd842b88
% canonical-lessons: ES-C670-indicar, ES-C670-indicacion, ES-C670-doblar, ES-C670-atravesar, ES-C670-desviarse

\chapter{Cómo llegar}
\label{ch:es-como-llegar}
\hlchaptermodality{670}

\begin{chapteropening}
\textbf{By the end of this chapter:} \emph{I can ask the way and follow directions: turn, go through and turn off.}

Complete the last lesson of chapter 670: I can ask the way and follow directions: turn, go through and turn off.
\end{chapteropening}

\section[indicar]{indicar — to show, to point out, to indicate}

\label{lesson:ES-C670-indicar}



\begin{quote}
Before the new one: say the Spanish for to point to, then the Spanish for located.
\end{quote}

\subsection*{You'll want to know: indicar}
\textbf{indicar} — \textquotedblleft{}to show, to point out, to indicate\textquotedblright{}.

\begin{quote}
¿Me puede indicar el camino a la estación?
\end{quote}

\begin{quote}
\emph{Can you show me the way to the station?}
\end{quote}

A polite way to ask for directions.

\begin{cousinweb}
\emph{indicar} is a learned borrowing from Latin \emph{indicāre} \textquotedblleft{}to point out, to show\textquotedblright{} (Lewis \& Short; EtymDB).
\end{cousinweb}

\begin{grammarlens}[title={What you've built}]
One more word for telling, explaining and describing at B1.
\end{grammarlens}

\subsection*{Your turn}
\begin{itemize}
\raggedright
  \item \emph{Say it:} \emph{indicar}
  \item \emph{Say it:} \emph{indicar}, once more
  \item \emph{Say it:} a sentence of your own with \emph{indicar}
  \item \emph{Recall:} say the Spanish for to point to, then the Spanish for located, then say \emph{indicar} again
\end{itemize}

\begin{tcolorbox}[breakable,colback=teal!4,colframe=teal!35!black,title={Before you move on}]
What does indicar mean? (\textquotedblleft{}To show, to point out, to indicate\textquotedblright{}.) Say it once more.
\end{tcolorbox}
\section[la indicación]{la indicación — the instruction, the direction (how to get somewhere)}

\label{lesson:ES-C670-indicacion}



\begin{quote}
Before the new one: say the Spanish for located, then the Spanish for to show.
\end{quote}

\subsection*{You'll want to know: la indicación}
\textbf{la indicación} — \textquotedblleft{}the instruction, the direction (how to get somewhere)\textquotedblright{}.

\begin{quote}
Sigue las indicaciones del hotel.
\end{quote}

\begin{quote}
\emph{Follow the hotel’s directions.}
\end{quote}

In the plural, \emph{las indicaciones} are directions.

\begin{cousinweb}
\emph{indicación} is a learned borrowing from Latin \emph{indicātiō} (EtymDB), from \emph{indicāre} \textquotedblleft{}to point out\textquotedblright{}.
\end{cousinweb}

\begin{grammarlens}[title={What you've built}]
One more word for telling, explaining and describing at B1.
\end{grammarlens}

\subsection*{Your turn}
\begin{itemize}
\raggedright
  \item \emph{Say it:} \emph{la indicación}
  \item \emph{Say it:} \emph{la indicación}, once more
  \item \emph{Say it:} a sentence of your own with \emph{la indicación}
  \item \emph{Recall:} say the Spanish for located, then the Spanish for to show, then say \emph{la indicación} again
\end{itemize}

\begin{tcolorbox}[breakable,colback=teal!4,colframe=teal!35!black,title={Before you move on}]
What does la indicación mean? (\textquotedblleft{}The instruction, the direction (how to get somewhere)\textquotedblright{}.) Say it once more.
\end{tcolorbox}
\section[doblar]{doblar — to turn (a corner); to fold; to dub}

\label{lesson:ES-C670-doblar}



\begin{quote}
Before the new one: say the Spanish for to show, then the Spanish for the instruction.
\end{quote}

\subsection*{You'll want to know: doblar}
\textbf{doblar} — \textquotedblleft{}to turn (a corner); to fold; to dub\textquotedblright{}.

\begin{quote}
Dobla a la izquierda en la plaza.
\end{quote}

\begin{quote}
\emph{Turn left at the square.}
\end{quote}

Very common for turning in Latin America; Spain also says \emph{girar}. \emph{Una película doblada} is a dubbed film.

\begin{cousinweb}
\emph{doblar} is inherited from Late Latin \emph{duplāre} \textquotedblleft{}to double\textquotedblright{} (Lewis \& Short; EtymDB), built on \emph{duplus} \textquotedblleft{}double\textquotedblright{}.
\end{cousinweb}

\begin{grammarlens}[title={What you've built}]
One more word for telling, explaining and describing at B1.
\end{grammarlens}

\subsection*{Your turn}
\begin{itemize}
\raggedright
  \item \emph{Say it:} \emph{doblar}
  \item \emph{Say it:} \emph{doblar}, once more
  \item \emph{Say it:} a sentence of your own with \emph{doblar}
  \item \emph{Recall:} say the Spanish for to show, then the Spanish for the instruction, then say \emph{doblar} again
\end{itemize}

\begin{tcolorbox}[breakable,colback=teal!4,colframe=teal!35!black,title={Before you move on}]
What does doblar mean? (\textquotedblleft{}To turn (a corner); to fold; to dub\textquotedblright{}.) Say it once more.
\end{tcolorbox}
\section[atravesar]{atravesar — to cross, to go through}

\label{lesson:ES-C670-atravesar}



\begin{quote}
Before the new one: say the Spanish for the instruction, then the Spanish for to turn.
\end{quote}

\subsection*{You'll want to know: atravesar}
\textbf{atravesar} — \textquotedblleft{}to cross, to go through\textquotedblright{}.

\begin{quote}
Atravesamos el parque a pie.
\end{quote}

\begin{quote}
\emph{We went through the park on foot.}
\end{quote}

The e becomes ie under stress: \emph{atravieso}, \emph{atraviesa}.

\begin{cousinweb}
\emph{atravesar} is built on \emph{través} \textquotedblleft{}crossways\textquotedblright{}, which is inherited from Latin \emph{trānsversus} \textquotedblleft{}lying across\textquotedblright{} (EtymDB), from \emph{trans} and \emph{vertere} \textquotedblleft{}to turn\textquotedblright{}.
\end{cousinweb}

\begin{grammarlens}[title={What you've built}]
One more word for telling, explaining and describing at B1.
\end{grammarlens}

\subsection*{Your turn}
\begin{itemize}
\raggedright
  \item \emph{Say it:} \emph{atravesar}
  \item \emph{Say it:} \emph{atravesar}, once more
  \item \emph{Say it:} a sentence of your own with \emph{atravesar}
  \item \emph{Recall:} say the Spanish for the instruction, then the Spanish for to turn, then say \emph{atravesar} again
\end{itemize}

\begin{tcolorbox}[breakable,colback=teal!4,colframe=teal!35!black,title={Before you move on}]
What does atravesar mean? (\textquotedblleft{}To cross, to go through\textquotedblright{}.) Say it once more.
\end{tcolorbox}
\section[desviarse]{desviarse — to turn off, to leave the road; to go off course}

\label{lesson:ES-C670-desviarse}



\begin{quote}
Before the new one: say the Spanish for to turn, then the Spanish for to cross.
\end{quote}

\subsection*{You'll want to know: desviarse}
\textbf{desviarse} — \textquotedblleft{}to turn off, to leave the road; to go off course\textquotedblright{}.

\begin{quote}
Desvíate en la próxima salida.
\end{quote}

\begin{quote}
\emph{Turn off at the next exit.}
\end{quote}

The i takes a written accent in some forms: \emph{te desvías}, \emph{se desvía}.

\begin{cousinweb}
\emph{desviarse} is built on \emph{desviar}, from Late Latin \emph{dēviāre} \textquotedblleft{}to turn aside, to stray\textquotedblright{} (EtymDB), built on \emph{via} \textquotedblleft{}road\textquotedblright{}.
\end{cousinweb}

\begin{grammarlens}[title={What you've built}]
One more word for telling, explaining and describing at B1.
\end{grammarlens}

\subsection*{Your turn}
\begin{itemize}
\raggedright
  \item \emph{Say it:} \emph{desviarse}
  \item \emph{Say it:} \emph{desviarse}, once more
  \item \emph{Say it:} a sentence of your own with \emph{desviarse}
  \item \emph{Recall:} say the Spanish for to turn, then the Spanish for to cross, then say \emph{desviarse} again
\end{itemize}

\begin{tcolorbox}[breakable,colback=teal!4,colframe=teal!35!black,title={Before you move on}]
What does desviarse mean? (\textquotedblleft{}To turn off, to leave the road; to go off course\textquotedblright{}.) Say it once more.
\end{tcolorbox}
